\chapter[Envelope Theory]{Envelope Theory: Luttinger Khon, Burt, and Ermoneit}
\index{envelope-function theory|textbf}
\index{Luttinger--Kohn model|textbf}
\index{Burt envelope-function theory|textbf}
\index{effective-mass model!Luttinger--Kohn}

\label{app:envelope-theory}

\chapterstatus{literature explanation and project interpretation. Luttinger and
Kohn (1955), Burt (1992, 1999), and Ermoneit et al. (2026) have been read in
full. This appendix explains
their constructions in the context of the present reduction problem; it does
not claim that these papers prove the project's modern Wannier, alignment,
impurity-extraction, or continuum-convergence statements.}

\section{Why these papers matter to the present problem}
\index{effective-mass equation!derivation}
\index{periodic parent}
\index{gentle perturbation}

The central problem of this monograph is to identify when a microscopic or
first-principles semiconductor operator can be replaced by a smaller lattice or
continuum operator without silently changing the state space, physical branch,
or intended observable. Luttinger and Kohn's 1955 paper addresses an early and
foundational version of that problem~\cite{luttingerKohn1955}. It begins with a
one-electron Hamiltonian
\begin{equation}
  \hat H
  =
  \hat H_0+\hat U,
  \label{eq:lk-parent-perturbation}
\end{equation}
where $\hat H_0$ is periodic and $\hat U$ is an additional perturbing potential.
The paper asks how the periodic problem can be replaced, near a selected band
edge, by an equation for a slowly varying envelope.

Burt's 1992 and 1999 papers extend this question from a gentle perturbation of
one bulk crystal to piecewise-periodic microstructures that may contain
atomically abrupt interfaces~\cite{burt1992,burt1999}. They distinguish an exact,
invertible envelope-function representation from the approximate local
differential equations conventionally called envelope or effective-mass models.
Their central methodological shift is to derive exact equations first and to
base subsequent approximations on the Fourier content of the envelope functions,
rather than assuming that the material composition or microscopic perturbation
itself varies slowly.

This is close to, but not identical with, the present project. Here the parent
may be a represented Kohn--Sham operator, the pristine and doped calculations
may be produced independently, the extracted impurity term may be nonlocal or
matrix valued, and a Wannier or other retained basis may intervene before a
continuum limit is considered. Luttinger and Kohn therefore provide a model for
how a controlled reduction can be organized, not a ready-made proof of the
entire modern pipeline.

\section{The periodic parent and the adapted basis}
\index{Bloch state!band-edge basis}
\index{state space!Luttinger--Kohn basis}

Let the Bloch eigenstates of the periodic parent satisfy
\begin{equation}
  \hat H_0\lvert\psi_{n\mathbf k}\rangle
  =
  \varepsilon_n(\mathbf k)
  \lvert\psi_{n\mathbf k}\rangle,
\end{equation}
with position-space form
\begin{equation}
  \psi_{n\mathbf k}(\mathbf r)
  =
  e^{i\mathbf k\cdot\mathbf r}
  u_{n\mathbf k}(\mathbf r),
\end{equation}
where $u_{n\mathbf k}$ is lattice periodic. For a simple band with its minimum
at $\mathbf k_0=\mathbf 0$, Luttinger and Kohn do not expand the perturbed state
directly in ordinary Bloch states or in Wannier functions. They introduce the
adapted family corresponding to their equation (II.4),
\begin{equation}
  \chi_{n\mathbf k}(\mathbf r)
  =
  e^{i\mathbf k\cdot\mathbf r}u_{n\mathbf 0}(\mathbf r).
  \label{eq:lk-adapted-basis}
\end{equation}
They show that this family is complete and orthonormal when the underlying
Bloch family is complete and orthonormal. The perturbed state can therefore be
written as
\begin{equation}
  \Psi(\mathbf r)
  =
  \sum_n\int_{\mathrm{BZ}}
  A_n(\mathbf k)\chi_{n\mathbf k}(\mathbf r)\,d\mathbf k.
  \label{eq:lk-state-expansion}
\end{equation}

This choice is important conceptually. The reduced equation is not obtained
only by differentiating a plotted band. A state-space representation is chosen
first, the microscopic state is expanded in that representation, and only then
are interband terms approximated or removed. In the language of this
monograph, equation~\eqref{eq:lk-adapted-basis} supplies an analytical
identification tied to one periodic parent. It is not a maximally localized
Wannier gauge, an independently fitted alignment map, or a numerical
pristine--doped basis match.

\section{The gentle-potential approximation}
\index{gentle perturbation!Fourier content}
\index{central-cell correction}

The main physical assumption is that the fractional change of $U(\mathbf r)$
over one unit cell is small. In reciprocal space, the matrix element of $U$
between adapted states contains Fourier components shifted by reciprocal
lattice vectors. Luttinger and Kohn retain the component with zero reciprocal
lattice vector when the perturbation is sufficiently smooth. In schematic
notation this gives
\begin{equation}
  \langle\chi_{n\mathbf k}\rvert
  \hat U
  \lvert\chi_{n'\mathbf k'}\rangle
  \approx
  \delta_{nn'}\,\widetilde U(\mathbf k-\mathbf k').
  \label{eq:lk-gentle-potential}
\end{equation}
The approximation suppresses large reciprocal-lattice momentum transfers and
makes the leading perturbation band diagonal.

The paper repeatedly connects this assumption to scale separation. If $a$ is a
lattice spacing and $a_i$ is the characteristic extent of an impurity state,
the neglected contributions in the simple-band derivation are estimated at
order $(a/a_i)^2$ under the stated assumptions. This is not a universal error
bound for an arbitrary first-principles impurity operator. It is an
order-of-magnitude estimate inside the paper's local-potential and slowly
varying-envelope setting.

Luttinger and Kohn also identify the central-cell limitation explicitly. If the
potential changes strongly inside the cell containing the impurity, the simple
envelope equation is not valid there. They argue that an envelope description
may remain valid sufficiently far from a short-range region. For the present
project, this distinction motivates keeping a short-range central-cell operator
separate from any proposed long-range screened continuum term rather than
forcing one model to describe both regimes without evidence.

\section{Removing interband coupling}
\index{canonical transformation!Luttinger--Kohn}
\index{interband coupling}
\index{band truncation}

In the adapted basis, the periodic Hamiltonian contains interband momentum
matrix elements proportional to the wave vector measured from the band edge.
Luttinger and Kohn treat those terms as small and introduce a canonical
transformation
\begin{equation}
  \mathbf A=e^{\mathbf S}\mathbf B
\end{equation}
chosen to eliminate first-order interband coupling. Expanding the transformed
Hamiltonian gives
\begin{equation}
  e^{-\mathbf S}\mathbf H e^{\mathbf S}
  =
  \mathbf H+[\mathbf H,\mathbf S]
  +\frac{1}{2}[[\mathbf H,\mathbf S],\mathbf S]+\cdots.
\end{equation}
The retained second-order terms reproduce the quadratic expansion of the
selected band energy. This is a perturbative decoupling argument, not exact
removal of the excluded bands for arbitrary energy, perturbation strength, or
wave vector.

The distinction matters for the current research program. Modern retained-band
selection, Wannier disentanglement, Feshbach downfolding, real-space
truncation, and model fitting are different operations. The 1955 canonical
transformation is evidence for the importance of accounting for excluded-band
coupling; it does not make those later operations synonymous.

\section{The simple-band envelope equation}
\index{envelope function!simple band}
\index{effective-mass tensor!band curvature}

After the smooth-potential and quadratic-band approximations, the transformed
coefficient for a simple band satisfies the momentum-space equation given as
(II.35) in the paper. Introducing the Fourier-transformed envelope
\begin{equation}
  F_n(\mathbf r)
  =
  \int_{\mathrm{BZ}}
  e^{i\mathbf k\cdot\mathbf r}B_n(\mathbf k)\,d\mathbf k
\end{equation}
leads first to a short-range convolution form and then, when its kernel can be
replaced by a delta distribution on the envelope scale, to the paper's equation
(II.38):
\begin{equation}
  \left[
    \varepsilon_n(-i\boldsymbol\nabla)
    +U(\mathbf r)
  \right]F_n(\mathbf r)
  =
  E F_n(\mathbf r).
  \label{eq:lk-envelope-operator}
\end{equation}
Here $\varepsilon_n$ is understood only through its quadratic expansion about
the selected extremum. Writing
\begin{equation}
  \varepsilon_n(\mathbf k_0+\mathbf q)
  =
  \varepsilon_n(\mathbf k_0)
  +
  \frac{\hbar^2}{2}
  \mathbf q^{\mathsf T}\mathbf m_n^{*-1}\mathbf q
  +O(\lVert\mathbf q\rVert^3)
\end{equation}
and subtracting the constant edge energy yields the familiar anisotropic
kinetic operator
\begin{equation}
  -\frac{\hbar^2}{2}
  \boldsymbol\nabla\mathbin{\cdot}
  \mathbf m_n^{*-1}
  \boldsymbol\nabla.
\end{equation}

The leading microscopic wavefunction is the product stated in the paper's
equation (II.40),
\begin{equation}
  \Psi(\mathbf r)
  \approx
  F_n(\mathbf r)\psi_{n\mathbf 0}(\mathbf r).
  \label{eq:lk-envelope-bloch-product}
\end{equation}
Equation~\eqref{eq:lk-envelope-bloch-product} displays the two-scale structure
that the present monograph seeks to recover: a slowly varying envelope and a
lattice-periodic microscopic factor. It is a leading approximation within a
single-parent construction, not an exact identity between independently
produced atomistic and continuum state spaces.

\section{Band minima away from the zone center and multiple valleys}
\index{multivalley model!Luttinger--Kohn}
\index{silicon!conduction valleys}
\index{valley!linear combination}

Section~III shifts the adapted basis to a minimum at
$\mathbf k_0\ne\mathbf 0$. The same derivation then uses band quantities
evaluated at that minimum, and the envelope equation contains the quadratic
expansion about $\mathbf k_0$. When symmetry produces several equivalent
minima, the leading state has the schematic form
\begin{equation}
  \Psi(\mathbf r)
  \approx
  \sum_{\nu=1}^{N_v}
  c_\nu F_\nu(\mathbf r)
  \psi_{n\mathbf k_\nu}(\mathbf r).
  \label{eq:lk-multivalley-state}
\end{equation}

The paper makes an important limitation explicit: the lowest-order
effective-mass equations do not determine the correct linear combinations of
degenerate valley solutions. Those combinations and their small splittings
require corrections beyond the leading theory. For substitutional phosphorus
in silicon, this prevents the existence of six equivalent conduction minima
from being mistaken for a completely specified donor model. Valley ordering,
phases, intervalley coupling, and the central-cell mechanism must still be
identified.

\section{Degenerate bands and coupled envelopes}
\index{degenerate band!coupled envelope equations}
\index{multiband model!Luttinger--Kohn}

For a band edge with an $r$-dimensional degeneracy, Section~IV replaces the
single envelope by components $F_i$. The quadratic band-edge information is a
matrix-valued form with coefficients $D_{ij}^{\alpha\beta}$, and the impurity
problem becomes a coupled system of the schematic form
\begin{equation}
  \sum_{j=1}^{r}
  \left[
    -\sum_{\alpha,\beta}
    D_{ij}^{\alpha\beta}
    \partial_\alpha\partial_\beta
    +U(\mathbf r)\delta_{ij}
  \right]F_j(\mathbf r)
  =
  E F_i(\mathbf r).
  \label{eq:lk-degenerate-envelope}
\end{equation}
The exact coefficients and sign conventions must be taken from the declared
band-edge basis and the paper's definitions. The structural conclusion is that
degeneracy changes the effective equation from one scalar PDE to a coupled
matrix differential equation.

This is directly relevant to the boron branch. A valence-band acceptor cannot
be justified by copying the scalar donor equation and changing the sign of the
charge. The retained valence manifold, its degeneracies, spin--orbit structure,
and allowed couplings determine the continuum state space.

\section{Spin--orbit coupling and the valence manifold}
\index{spin--orbit coupling!Luttinger--Kohn}
\index{acceptor!coupled valence manifold}

Section~V incorporates spin--orbit coupling into the band-edge matrix elements.
For a nondegenerate band, the formal structure of the effective-mass equation
survives with modified momentum matrix elements and the spin--orbit-corrected
band dispersion. For a degenerate valence edge, spin--orbit coupling partially
lifts the degeneracy and yields coupled multiplets.

The paper emphasizes a scale criterion: a single spin--orbit-split band can be
treated independently only when the energies induced by the impurity or
magnetic field are small compared with the separation to the neighboring
spin--orbit band. In its discussion of silicon acceptors, it warns that the
binding-energy and spin--orbit scales can be comparable, so the simple
single-band effective-mass theory need not be valid. The paper consequently
constructs a coupled six-component description before considering decoupled
limits.

For this monograph, the lesson is methodological rather than numerical. Final
boron claims require a spinor parent and a declared coupled valence-manifold
continuum model. The historical numerical scale estimates in the 1955 paper are
not adopted as calculated or current reference values for this project without
separate verification.

\section{Magnetic fields}
\index{magnetic field!effective-mass equation}

Although magnetic-field calculations are outside the initial project scope,
the paper's treatment is conceptually useful. It shows that the adapted-basis
and canonical-transformation method extends to an external homogeneous field.
At quadratic order, the effective equation follows the replacement
\begin{equation}
  \mathbf k
  \longmapsto
  -i\boldsymbol\nabla-\frac{e}{\hbar c}\mathbf A,
\end{equation}
with noncommuting products symmetrized. This part of the paper demonstrates
that an effective operator must preserve the algebra required by the applied
field; matching only the zero-field dispersion would not establish that
property.

\section{Burt's exact envelope-function representation}
\index{Burt envelope-function theory!exact representation}
\index{envelope function!band-limited definition}
\index{Brillouin-zone restriction}

The equations in this and the following sections are explicit extractions from
Burt's papers, rewritten only to match the notation of this monograph; the
source equation numbers are stated at each step. Burt retains the
Luttinger--Kohn form of the expansion but changes its scope and
use~\cite{burt1992,burt1999}. Choose one complete periodic basis
$\{U_n(\mathbf r)\}$ for the entire structure and write
\begin{equation}
  \Psi(\mathbf r)
  =
  \sum_n F_n(\mathbf r)U_n(\mathbf r),
  \label{eq:burt-envelope-expansion}
\end{equation}
with every envelope $F_n$ defined to contain only plane waves whose wave vectors
lie in one chosen primitive reciprocal cell, usually a Brillouin zone. This is
Burt's 1992 equation (3.1) and 1999 equation (2.4). If
\begin{equation}
  U_n(\mathbf r)
  =
  \sum_{\mathbf G}U_{n\mathbf G}e^{i\mathbf G\cdot\mathbf r},
\end{equation}
then the explicit envelope extraction in Burt's 1999 equation (2.5) is
\begin{equation}
  F_n(\mathbf r)
  =
  \sum_{\mathbf k\in\mathrm{BZ}}
  \sum_{\mathbf G}
  U_{n\mathbf G}^{*}
  \widetilde\Psi(\mathbf k+\mathbf G)
  e^{i\mathbf k\cdot\mathbf r}.
  \label{eq:burt-explicit-envelope-coefficient}
\end{equation}
For an orthonormal periodic basis, this is a unitary transformation at each
in-zone wave vector. The expansion is therefore complete, unique, and
invertible before it is truncated.

This definition differs from a conventional regional ansatz. In a conventional
heterostructure treatment, one may choose different bulk band-edge functions in
different layers and then match trial solutions at an interface. Burt instead
uses the same periodic basis throughout, even where those functions do not
diagonalize the local bulk Hamiltonian. The material change appears in the
Hamiltonian matrix, including its off-diagonal entries, rather than through an
unrecorded change of basis.

The band limit has two consequences that must be kept together. First, the exact
$F_n$ are smooth and continuous, with continuous derivatives, for a sufficiently
well-behaved microscopic field. Second, one envelope by itself need not resemble
the original wavefunction. Completeness belongs to the full sum in
\eqref{eq:burt-envelope-expansion}. Burt's 1999 delta-function example makes this
point sharply: each individual band-limited envelope is broad and oscillatory,
but the infinite collection reconstructs the delta distribution exactly.

\section{Exact equations before reduction}
\index{envelope-function equation!exact nonlocal form}
\index{nonlocal operator!Burt envelope representation}

Inserting equation~\eqref{eq:burt-envelope-expansion} into the microscopic
Schr\"odinger equation and applying the inverse basis transformation gives
Burt's exact 1992 equation (3.8), which can be written
\begin{equation}
  -\frac{\hbar^2}{2m}\nabla^2F_n
  -\frac{i\hbar}{m}\sum_{n'}
    \mathbf p_{nn'}\mathbin{\cdot}\boldsymbol\nabla F_{n'}
  +\sum_{n'}\int
    H_{nn'}(\mathbf r,\mathbf r')F_{n'}(\mathbf r')
    \,d\mathbf r'
  =
  E F_n(\mathbf r).
  \label{eq:burt-exact-envelope}
\end{equation}
Here the periodic-basis kinetic and momentum matrix elements are contained in
$H_{nn'}$ and $\mathbf p_{nn'}$ according to the paper's definitions, and the
restricted delta kernel appearing in its equations (3.9)--(3.10) is
\begin{equation}
  \Delta_{\mathrm{BZ}}(\mathbf r-\mathbf r')
  =
  \frac{1}{\Omega}
  \sum_{\mathbf k\in\mathrm{BZ}}
  e^{i\mathbf k\cdot(\mathbf r-\mathbf r')}.
  \label{eq:burt-restricted-delta}
\end{equation}
The remaining kernel contains the transformed potential operator. It can be
nonlocal even when the microscopic potential is local because the envelope
Fourier domain is restricted to a Brillouin zone. Burt's alternative 1992
derivation exhibits equation~\eqref{eq:burt-exact-envelope} as a unitary change
of representation from a plane-wave Hamiltonian. A truncated plane-wave
Hamiltonian and its correspondingly transformed envelope equations therefore
have identical eigenvalues before further approximation.

The word \emph{exact} in this context refers to the representation transformation
and the complete set of envelope equations. It does not mean that one dominant
envelope, a local effective-mass equation, a finite band set, or a
piecewise-constant coefficient model is exact. Those simplifications enter only
later and must be assessed separately.

This distinction is particularly relevant to the present project. A nonlocal
pseudopotential or represented impurity operator is not excluded by the exact
envelope transformation. But obtaining a local scalar continuum potential from
such an operator is an additional model reduction, not a consequence of naming
$F_n$ an envelope.

\section{Global slow variation rather than a gentle structure}
\index{globally slowly varying envelope}
\index{abrupt interface!envelope theory}

Luttinger and Kohn focus on a perturbing potential that is gentle on the lattice
scale. Burt's microstructure theory changes the approximation question. The
microscopic composition may change abruptly; what matters for the simplified
envelope equations is whether the relevant envelopes are \emph{globally slowly
varying}: their Fourier weight is concentrated near the selected in-zone wave
vectors~\cite{burt1992,burt1999}.

This condition does not require every pointwise derivative to be small.
Approximate envelopes may display a localized rapid change, kink, or interface
feature while remaining dominated globally by small wave vectors. Conversely,
a smooth-looking plotted function is not sufficient evidence: the Fourier
support and the omitted high-wave-vector contributions must be checked.

Under the globally-slow condition, the nonlocal kernel in
\eqref{eq:burt-exact-envelope} may be approximated by a local or short-range
operator. Burt's 1992 equation (5.1) then has the explicit local multiband form
\begin{equation}
  -\frac{\hbar^2}{2m}\nabla^2F_n
  -\frac{i\hbar}{m}\sum_{n'}
    \mathbf p_{nn'}\mathbin{\cdot}\boldsymbol\nabla F_{n'}
  +\sum_{n'}H_{nn'}(\mathbf r)F_{n'}
  =
  E F_n.
  \label{eq:burt-local-multiband}
\end{equation} Burt separates bulk-profile terms, determined by the constituent
bulk Hamiltonians and geometry, from interface terms that depend on how the
microscopic potential changes. For the model analyzed in 1992, interface terms
carry large reciprocal-wave-vector denominators and can be small for the
low-lying state, but the paper does not claim that all interface terms are
universally negligible. The appropriate test is made after the target state and
operator have been identified.

For an impurity problem, the analogous project question is not simply whether
the impurity atom is localized. It is whether the aligned microscopic state has
envelope Fourier weight concentrated in the retained sectors and whether the
terms discarded in localizing the envelope kernel are small on the declared
comparison domain.

\section{Dominant envelopes and effective-mass reduction}
\index{dominant envelope}
\index{folding down!Burt envelope theory}
\index{effective-mass operator!position dependent}

Burt obtains effective-mass equations by partitioning the envelopes into a
small dominant set and a remote set, then eliminating the remote components
perturbatively. This is a folding-down step performed after the exact envelope
representation has been established. For a dominant scalar conduction envelope,
the resulting operator has the schematic form
\begin{equation}
  -\frac{\hbar^2}{2}
  \boldsymbol\nabla\mathbin{\cdot}
  \mathbf m^{-1}(E,\mathbf r)
  \boldsymbol\nabla
  +E_c^{\mathrm{eff}}(\mathbf r),
  \label{eq:burt-position-dependent-mass}
\end{equation}
where the energy and position dependence arise from the remote-band elimination
and the transformed local Hamiltonian. This is the operator structure extracted
from Burt's 1992 equations (6.9)--(6.10). The coefficients are derived rather
than selected by choosing one of several formally Hermitian orderings.

The 1992 derivation keeps changes in the constituent zone-center eigenstates by
retaining off-diagonal matrix elements of the zone-center Hamiltonian in one
fixed basis. This explains why an effective-mass equation can remain useful even
when the well and barrier Bloch functions are not approximately identical. It
also identifies the approximation: interface terms, higher derivatives, remote
envelopes, or nonlocal kernel contributions are neglected according to the
chosen order.

Energy-dependent effective masses can arise naturally in this elimination. Burt
distinguishes an energy-eigenvalue equation with energy-dependent coefficients
from a time-dependent evolution equation and argues that the former is not
forbidden by the superposition principle. Such a mass remains an approximation
whose usefulness depends on the coupling structure and the neglected terms; it
is not a universal prescription for every nonparabolic band.

\section{Operator ordering and whole-matrix Hermiticity}
\index{operator ordering!Burt envelope theory}
\index{Hermiticity!whole multiband Hamiltonian}

Burt's 1999 tutorial derivation emphasizes that replacing constant bulk
parameters by functions of position and symmetrizing every matrix element is not
a reliable way to construct a multiband Hamiltonian~\cite{burt1999}. The correct
order of derivatives and coefficients follows from the order of momentum matrix
elements in the remote-band elimination.

Burt's 1999 equations (4.4d)--(4.6) show the extraction directly. In the
four-band example, eliminating the small conduction envelope gives
\begin{equation}
  F_S
  \approx
  \frac{i\hbar p}{mE_g}
  \left(
    \partial_xF_X+\partial_yF_Y+\partial_zF_Z
  \right),
  \label{eq:burt-eliminated-conduction-envelope}
\end{equation}
and the $X$-component equation becomes
\begin{equation}
  \left(E_V-E-\frac{\hbar^2}{2m}\nabla^2\right)F_X
  +
  \partial_x\!\left[
    \frac{P^2}{E_g}
    \left(
      \partial_xF_X+\partial_yF_Y+\partial_zF_Z
    \right)
  \right]
  =0,
  \label{eq:burt-valence-component}
\end{equation}
with cyclic equations for $F_Y$ and $F_Z$. In $\mathbf k=-i\boldsymbol\nabla$
notation, Burt's equation (4.11) places an off-diagonal coefficient as
$-k_i b(\mathbf r)k_j$, with the order inherited from the two momentum matrix
elements rather than from entrywise symmetrization.

For a coupled valence system, an individual off-diagonal differential operator
need not be Hermitian by itself. Hermiticity is a property of the complete
Hamiltonian matrix: the adjoint partner occurs in the transposed block. Different
remote-band contributions to what is one bulk Luttinger parameter can acquire
different operator orderings in an inhomogeneous structure. Consequently, a
single fitted bulk parameter may need to be resolved into its microscopic
contributions before an interface operator is defined.

The project analogue is immediate. Hermiticity of a complete represented
impurity or continuum operator does not require every labeled physical component
to be independently Hermitian before assembly. Conversely, entrywise
symmetrization cannot substitute for a derivation from the chosen state-space
map and eliminated sectors.

\section{Interfaces, boundary conditions, and current}
\index{boundary condition!Burt envelope theory}
\index{probability current!interface}
\index{effective-mass kink}

For the exact band-limited envelopes, continuity is built into the definition;
it is not imposed by matching unrelated regional expansions. Integrating the
local multiband equation gives the explicit interface quantity extracted from
Burt's 1992 equations (7.1)--(7.2):
\begin{equation}
  -\frac{\hbar^2}{2m}\,\partial_{\perp}F_n
  -\frac{i\hbar}{m}\sum_{n'}p^{\perp}_{nn'}F_{n'}
  \quad\text{is continuous across the interface},
  \label{eq:burt-interface-flux}
\end{equation}
where $p^{\perp}$ is the momentum component normal to the interface. After a local
multiband approximation is made, retaining the full second-order system gives
continuity conditions for the envelopes and their derivatives in the setting
analyzed by Burt. If remote envelopes or second derivatives are then eliminated,
the reduced scalar equation~\eqref{eq:burt-position-dependent-mass} gives the
familiar condition that the normal component of
\begin{equation}
  \mathbf m^{-1}(E,\mathbf r)\boldsymbol\nabla F
\end{equation}
is continuous across an abrupt effective-mass step.

The corresponding kink in the approximate scalar envelope does not imply that
the exact microscopic wavefunction has a literal derivative discontinuity.
Burt's one-dimensional model shows a rapid but smooth change of derivative in
the exact wavefunction and exact principal envelope; the piecewise-constant
mass model replaces this soft interface region by a kink. The same derived
boundary condition is consistent with the probability current of the reduced
equation.

Thus boundary conditions belong to the approximation used. They must not be
borrowed from one scalar effective-mass equation and applied unchanged to a
multiband, nonlocal, or differently ordered operator.

\section{Luttinger--Kohn and Wannier--Slater envelopes}
\index{Wannier--Slater envelope}
\index{Luttinger--Kohn model!comparison with Wannier--Slater}

Burt distinguishes two envelope representations. In the Luttinger--Kohn form,
the periodic basis is independent of the in-zone wave vector:
\begin{equation}
  \Psi(\mathbf r)
  =
  \sum_{n\mathbf k}
  F_n(\mathbf k)e^{i\mathbf k\cdot\mathbf r}U_n(\mathbf r).
\end{equation}
In the Wannier--Slater form, the periodic basis is wave-vector dependent:
\begin{equation}
  \Psi(\mathbf r)
  =
  \sum_{n\mathbf k}
  F_n^{\mathrm{WS}}(\mathbf k)
  e^{i\mathbf k\cdot\mathbf r}u_{n\mathbf k}(\mathbf r).
\end{equation}
These are Burt's 1992 equations (9.1)--(9.2). The latter envelopes are related
to amplitudes in Wannier orbitals of a reference crystal by the paper's equation
(9.10),
\begin{equation}
  F_n^{\mathrm{WS}}(\mathbf r)
  =
  V_c^{1/2}\sum_{\mathbf R}
  \Delta_{\mathrm{BZ}}(\mathbf r-\mathbf R)
  \langle n\mathbf R\mid\Psi\rangle.
  \label{eq:burt-wannier-slater-amplitudes}
\end{equation}
Exact equations can be derived for either representation. In Burt's 1992
analysis, the fixed-basis Luttinger--Kohn route gives a less cumbersome
systematic effective-mass derivation for abrupt and substantial material changes,
whereas the Wannier--Slater perturbation expansion generates additional interface
terms and becomes more involved at consistent higher order.

Neither representation should be identified automatically with the modern
construction of maximally localized Wannier functions. A Wannier90 calculation
selects a retained subspace and gauge numerically; Burt's comparison concerns two
analytical definitions of envelope coefficients relative to a reference periodic
Hamiltonian. Any connection to the project's Wannier operator must therefore be
made by an explicit map.

\section{Out-of-zone solutions and representation completeness}
\index{out-of-zone solution}
\index{envelope function!representation versus trial solution}

The 1999 paper separates exact representation envelopes $F_n$ from conventional
trial envelopes that solve approximate local differential equations. The latter
are not restricted to in-zone Fourier components. When one solves a derived
multiband local system in piecewise-periodic regions, roots with wave vectors
outside the reference Brillouin zone can appear.

The approximate one-dimensional system used for the explicit discussion is
Burt's 1999 equation (5.11),
\begin{equation}
  -\frac{\hbar^2}{2m}\frac{d^2F_n}{dx^2}
  -\frac{i\hbar}{m}\sum_{n'}p_{nn'}\frac{dF_{n'}}{dx}
  +\sum_{n'}H_{nn'}(x)F_{n'}
  =E'F_n.
  \label{eq:burt-out-of-zone-system}
\end{equation}
Burt argues that its out-of-zone solutions cannot be rejected solely because
of their wave vector. In the extended local equation they can be replicas needed
to satisfy the full set of interface conditions and reconstruct the microscopic
state. In the one-dimensional example, using all such solutions gives an
accurate wavefunction, and a reconstruction from one displaced reciprocal zone
also succeeds. This does not make every out-of-zone root physically trustworthy:
use of the extended equation requires the neglected interzone coupling and
large-denominator terms to be controlled. The derivation and reconstruction,
not the label ``out of zone,'' determine admissibility.

For numerical reduction in this project, a large-wave-vector component should
therefore be diagnosed rather than deleted by convention. It may indicate a
necessary reconstruction component, a failure of the globally-slow assumption,
an inadequate retained basis, or an artifact of extrapolating an approximate
Hamiltonian beyond its validated domain.

\section{Observables are not determined by the dominant envelope alone}
\index{observable!envelope reconstruction}
\index{dipole matrix element!envelope theory}

Burt's 1999 discussion of interband dipole matrix elements supplies a broader
warning for model reduction. Even when one envelope component dominates the
norm and the product $F_nU_n$ closely approximates the wavefunction pointwise,
that term alone need not approximate a matrix element of the position operator.
Small envelope components and zone-boundary terms can determine the observable.
For the one-dimensional dominant-component contribution, Burt's 1999 equation
(6.12) isolates the zone-edge correction,
\begin{align}
  I_{cv}
  ={}&\pi i\,
  \widetilde F_c(\pi/a)^{*}\widetilde F_v(-\pi/a)
  \sum_G U_{cG}^{*}U_{v,G+2\pi/a}
  \notag\\
  &-\pi i\,
  \widetilde F_c(-\pi/a)^{*}\widetilde F_v(\pi/a)
  \sum_G U_{c,G+2\pi/a}^{*}U_{vG}.
  \label{eq:burt-dipole-zone-edge}
\end{align}
This term is often small for globally slow envelopes, but it is not identically
zero merely because the periodic basis functions belong to different bands.

The lesson for the present work is that state or band reconstruction error does
not automatically bound every observable error. Acceptance of a reduced impurity
model must identify the observable, transform its operator consistently into the
retained envelope space, and test whether discarded components enter at leading
order for that observable.

\section{Ermoneit's exact multivalley specialization}
\index{Ermoneit multivalley theory}
\index{valley-sector projector}
\index{energy-reference invariance}

Ermoneit et al. specialize Burt--Foreman envelope theory to a decomposition of
the Brillouin zone into disjoint valley sectors and apply it to valley splitting
in strained Si/SiGe nanostructures~\cite{ermoneit2026}. The word \emph{exact}
again refers to the envelope representation and its transformed equation. Their
microscopic model still assumes a periodic silicon potential $V$ and represents
the heterostructure, gate, and electric-field contributions through a chosen
mesoscopic potential $U$; exactness is not a claim that this parent potential is
an atomistically exact description of Si/SiGe.

Let $S_v=S(\mathbf k_v)$ denote the non-overlapping sector assigned to valley
center $\mathbf k_v$. Their equations (2), (5), and (6) read
\begin{align}
  \Psi_\alpha(\mathbf r)
  &=
  \sum_{n,v}u_{n\mathbf k_v}(\mathbf r)
  F_{nv\alpha}(\mathbf r),
  \label{eq:ermoneit-valley-expansion}\\
  F_{nv\alpha}(\mathbf r)
  &=
  \sum_{\mathbf K\in S_v}
  e^{i\mathbf K\cdot\mathbf r}
  \sum_{\mathbf G}c_{n\mathbf k_v}(\mathbf G)^{*}
  \widetilde\Psi_\alpha(\mathbf K+\mathbf G),
  \label{eq:ermoneit-sector-envelope}\\
  f_{nv\alpha}(\mathbf r)
  &=
  e^{-i\mathbf k_v\cdot\mathbf r}F_{nv\alpha}(\mathbf r).
  \label{eq:ermoneit-slow-envelope}
\end{align}
Here $V_c$ is the finite crystal or computational normalization volume and
$\Omega_p$ is one primitive-cell volume. The source uses the discrete Fourier
convention
\begin{equation}
  \Psi_\alpha(\mathbf r)
  =
  \sum_{\mathbf q}\widetilde\Psi_\alpha(\mathbf q)
  e^{i\mathbf q\cdot\mathbf r},
  \qquad
  \frac{1}{\Omega_p}\int_{\Omega_p}
  u_{n\mathbf k}(\mathbf r)^{*}u_{n'\mathbf k}(\mathbf r)\,d^3r
  =\delta_{nn'}.
  \label{eq:ermoneit-source-normalization}
\end{equation}
Equation (9a) correspondingly gives
\begin{equation}
  \sum_{n,v}\int_{V_c}
  F_{nv\alpha}(\mathbf r)^{*}F_{nv\alpha'}(\mathbf r)\,d^3r
  =\delta_{\alpha\alpha'}.
  \label{eq:ermoneit-envelope-normalization}
\end{equation}
Exact completeness requires all bands and a collection of valley sectors whose
disjoint union is the full first Brillouin zone. A later restriction to selected
bands or valleys is a reduction, not part of the exact completeness statement.
Unlike an informal sum of valley-centered trial functions, the disjoint support
condition $\mathbf K\in S_v$ makes this decomposition unique. Equations
(9b)--(11) make the associated projection explicit through
\begin{equation}
  \Delta_v(\mathbf r-\mathbf r')
  =
  \frac{1}{V_c}
  \sum_{\mathbf K\in S_v}
  e^{i\mathbf K\cdot(\mathbf r-\mathbf r')},
  \qquad
  \int_{V_c}\!\Delta_v(\mathbf r-\mathbf r')
  F_{n v'\alpha}(\mathbf r')\,d^3r'
  =
  \delta_{vv'}F_{nv\alpha}(\mathbf r).
  \label{eq:ermoneit-sector-projector}
\end{equation}
Thus $\Delta_v$ is a sector-projected delta kernel, not the unrestricted Dirac
delta distribution.

\subsection{The project--multiply--project potential kernel}
\index{nonlocal operator!Ermoneit multivalley theory}

The source's equation (14) is the concrete nonlocal potential operator
\begin{equation}
  U^{nn'}_{vv'}(\mathbf r,\mathbf r')
  =
  \int_{V_c}\!d^3x\,
  \Delta_v(\mathbf r-\mathbf x)
  u_{n\mathbf k_v}(\mathbf x)^{*}U(\mathbf x)
  u_{n'\mathbf k_{v'}}(\mathbf x)
  \Delta_{v'}(\mathbf x-\mathbf r').
  \label{eq:ermoneit-project-multiply-project}
\end{equation}
Its structure is literally sector projection, multiplication by the microscopic
Bloch-factor matrix element of $U$, and projection into the receiving sector.
The transformed slowly varying kernel is
\begin{equation}
  \mathsf u^{nn'}_{vv'}(\mathbf r,\mathbf r')
  =
  e^{-i(\mathbf k_v\cdot\mathbf r-
         \mathbf k_{v'}\cdot\mathbf r')}
  U^{nn'}_{vv'}(\mathbf r,\mathbf r'),
  \label{eq:ermoneit-slow-kernel}
\end{equation}
which is their equation (16). After perturbative interband decoupling, but
without discarding intervalley coupling, their single-conduction-band equation
(19) becomes
\begin{equation}
  E_\alpha f_{v\alpha}(\mathbf r)
  =
  \left[-\frac{\hbar^2}{2}
    \boldsymbol\nabla\mathbin{\cdot}\mathbf m_{cv}^{-1}
    \boldsymbol\nabla+E_{cv}\right]f_{v\alpha}(\mathbf r)
  +
  \sum_{v'}\int_{V_c}\!d^3r'\,
  \mathsf u_{vv'}(\mathbf r,\mathbf r')
  f_{v'\alpha}(\mathbf r').
  \label{eq:ermoneit-single-band-nonlocal}
\end{equation}
This equation is local in the effective kinetic term and nonlocal in the
valley-projected potential. Interband elimination has occurred; intervalley
coupling has not.

\subsection{Energy-reference invariance}
\index{energy zero!intervalley invariance}

A constant energy-reference shift must change the diagonal energy but not the
physical intervalley coupling. Ermoneit et al. prove this at the kernel level.
Appendix~B of the paper first derives the band-indexed identity
\begin{equation}
  U^{nn'}_{vv'}(\mathbf r,\mathbf r')
  \xrightarrow{\,U\mapsto U+U_0\,}
  U^{nn'}_{vv'}(\mathbf r,\mathbf r')
  +U_0\,\delta_{nn'}\delta_{vv'}
  \Delta_v(\mathbf r-\mathbf r').
  \label{eq:ermoneit-reference-shift-uppercase}
\end{equation}
After the phase transformation~\eqref{eq:ermoneit-slow-kernel}, their equation
(26) is the corresponding single-band slowly varying kernel identity
\begin{equation}
  \mathsf u_{vv'}(\mathbf r,\mathbf r')
  \xrightarrow{\,U\mapsto U+U_0\,}
  \mathsf u_{vv'}(\mathbf r,\mathbf r')
  +U_0\,\delta_{vv'}
  e^{-i\mathbf k_v\cdot(\mathbf r-\mathbf r')}
  \Delta_v(\mathbf r-\mathbf r').
  \label{eq:ermoneit-reference-shift}
\end{equation}
The phase-shifted projector acts as the identity on the corresponding slowly
varying envelope. Because the sectors are disjoint, the additional constant is
exactly diagonal in the valley label. Consequently,
\begin{equation}
  \hat H_0\mapsto\hat H_0+U_0 I,
  \qquad
  \hat H_1\mapsto\hat H_1,
  \label{eq:ermoneit-intravalley-intervalley-shift}
\end{equation}
where $\hat H_0$ and $\hat H_1$ are the intra- and intervalley blocks of their
equations (22)--(24).

Replacing $\Delta_v$ by an unrestricted Dirac delta produces the conventional
local model (21), but also removes the sector constraint. Under the same shift,
the paper obtains the additional local off-diagonal contribution
\begin{equation}
  \hat H_{1}^{\mathrm{loc}}
  \mapsto
  \hat H_{1}^{\mathrm{loc}}
  +U_0\sum_{v'\ne v}
  e^{-i(\mathbf k_v-\mathbf k_{v'})\cdot\mathbf r}
  u_{\mathbf k_v}(\mathbf r)^{*}
  u_{\mathbf k_{v'}}(\mathbf r).
  \label{eq:ermoneit-local-reference-artifact}
\end{equation}
This is not a physical valley-coupling response to a changed structure; it is an
artifact of extending local envelopes beyond their defining valley sectors.
For this problem, invariance under a constant energy shift is therefore a sharp
operator-level diagnostic of the local approximation.

\subsection{Intervalley coupling and valley splitting}
\index{valley splitting!nonlocal coupling}

For the two low-energy valleys at
$\mathbf k_0^{\pm}=(0,0,\pm k_0)^{\mathsf T}$, the unperturbed single-valley
operators satisfy $\hat H_0^{-}=\hat H_0^{+*}$. The zeroth-order modes are
chosen so that
\begin{equation}
  f_-^{(0)}=f_+^{(0)*},
  \qquad
  \int_{V_c}|f_+^{(0)}(\mathbf r)|^2\,d^3r
  =
  \int_{V_c}|f_-^{(0)}(\mathbf r)|^2\,d^3r
  =1.
  \label{eq:ermoneit-valley-mode-normalization}
\end{equation}
Degenerate perturbation theory reduces the first-order problem to
\begin{equation}
  \begin{pmatrix}0&\Delta\\\Delta^{*}&0\end{pmatrix}
  \begin{pmatrix}\eta_+\\\eta_-\end{pmatrix}
  =
  E_\alpha^{(1)}
  \begin{pmatrix}\eta_+\\\eta_-\end{pmatrix},
  \qquad
  E_{\mathrm{VS}}=2|\Delta|.
  \label{eq:ermoneit-valley-splitting-matrix}
\end{equation}
The first-order intervalley coupling evaluated with the exact sector-projected
nonlocal kernel in their equation (30) is
\begin{equation}
  \Delta
  =
  \int_{V_c}\!d^3r\int_{V_c}\!d^3r'\,
  f_{+}^{(0)}(\mathbf r)^{*}
  \mathsf u_{+-}(\mathbf r,\mathbf r')
  f_{-}^{(0)}(\mathbf r').
  \label{eq:ermoneit-nonlocal-intervalley-coupling}
\end{equation}
An exact algebraic rewriting of this first-order matrix element, using the
convolution property of the sector projectors, gives their equation (31),
\begin{equation}
  \Delta
  =
  \int_{V_c}\!d^3r\,
  e^{-2ik_0z}
  f_{+}^{(0)}(\mathbf r)^{*}
  u_{+}(\mathbf r)^{*}U(\mathbf r)
  u_{-}(\mathbf r)f_{-}^{(0)}(\mathbf r).
  \label{eq:ermoneit-local-looking-exact-coupling}
\end{equation}
Although this last expression looks like the conventional local-EFT formula,
its envelopes remain sector projected. The support restriction, rather than the
visual form of the final integral, is what preserves energy-reference
invariance.

\subsection{Spectrally filtered local approximation}
\index{spectral filtering!valley envelope}

The paper also gives an intermediate approximation that retains the inexpensive
local single-valley solve and then restores the valley-sector support. For the
one-dimensional reduction, define
\begin{equation}
  \Delta_v^{(1\mathrm D)}(z-z')
  =
  \frac{1}{L_z}\sum_{K_z\in S_v^{(z)}}e^{iK_z(z-z')},
  \label{eq:ermoneit-one-dimensional-projector}
\end{equation}
where $L_z$ is the normalization length and $S_v^{(z)}$ is the relevant
one-dimensional valley-sector interval. The paper's equation (39) is then
\begin{equation}
  \widetilde f_{v\alpha}^{(0)}(z)
  =
  \int dz'\,
  \Delta_v^{(1\mathrm D)}(z-z')e^{-ik_v(z-z')}
  f_{\mathrm{loc},\alpha}^{(0)}(z'),
  \label{eq:ermoneit-filtered-envelope}
\end{equation}
followed by normalization. Substitution of these projected envelopes into the
coupling integral gives equation (40),
\begin{equation}
  \Delta_{\mathrm{loc}}^{\mathrm{filtered}}
  =
  \sum_{n\in\mathbb Z}C_n^{(2)}
  \int dz\,
  e^{-i(2k_0+nG_{0z})z}
  \widetilde f_{+}^{(0)}(z)^{*}U_{\perp}(z)
  \widetilde f_{-}^{(0)}(z).
  \label{eq:ermoneit-filtered-coupling}
\end{equation}
By construction, this approximation is invariant under $U\mapsto U+U_0$.
It is not identical to the nonlocal theory. The authors' reported numerical
comparisons show close agreement for several quantum-well and wiggle-well cases,
while their sharp Ge-spike example shows that filtering can overestimate the
nonlocal result. These are literature-reported calculations for the stated
Si/SiGe model, not calculated results or scientific validation for the present
silicon-dopant project.

\section{Mapping the papers onto the project reduction chain}
\index{reduction!Luttinger--Kohn mapping}
\index{alignment!absent by construction in Luttinger--Kohn}

The relation between the historical construction and the present project can
be summarized as
\begin{equation}
  \begin{aligned}
  \text{periodic parent }\hat H_0
  &\longrightarrow
  \text{band-edge adapted basis}\\
  &\longrightarrow
  \text{perturbative interband decoupling}\\
  &\longrightarrow
  \text{envelope or coupled-envelope operator}.
  \end{aligned}
  \label{eq:lk-reduction-chain}
\end{equation}
Burt inserts an exact representation layer before any local reduction:
\begin{equation}
  \begin{aligned}
  \text{piecewise-periodic microscopic operator}
  &\longrightarrow
  \text{one fixed periodic basis}\\
  &\longrightarrow
  \text{exact band-limited nonlocal envelopes}\\
  &\longrightarrow
  \text{globally-slow Fourier diagnosis}\\
  &\longrightarrow
  \text{local multiband equation}\\
  &\longrightarrow
  \text{dominant-envelope effective operator}.
  \end{aligned}
  \label{eq:burt-reduction-chain}
\end{equation}
Ermoneit et al. then make the multivalley support constraint and one diagnostic
of its preservation explicit:
\begin{equation}
  \begin{aligned}
  \text{disjoint valley sectors}
  &\longrightarrow
  \text{sector-projected delta kernels}\\
  &\longrightarrow
  \text{project--multiply--project potential}\\
  &\longrightarrow
  \text{nonlocal intervalley coupling}\\
  &\longrightarrow
  \text{energy-reference-invariance test}\\
  &\longrightarrow
  \text{optional filtered-local approximation}.
  \end{aligned}
  \label{eq:ermoneit-reduction-chain}
\end{equation}
The corresponding modern project chain is longer:
\begin{equation}
  \begin{aligned}
  \text{pristine and doped KS parents}
  &\longrightarrow
  \text{compatible retained subspaces}\\
  &\longrightarrow
  \text{gauge and energy alignment}\\
  &\longrightarrow
  \text{aligned impurity operator}\\
  &\longrightarrow
  \text{lattice reduction}\\
  &\longrightarrow
  \text{continuum embedding and comparison}.
  \end{aligned}
  \label{eq:project-reduction-chain-lk}
\end{equation}

Several differences must remain explicit.

\begin{enumerate}
  \item \textbf{One microscopic problem versus independently produced
    calculations.} Luttinger and Kohn write $\hat H=\hat H_0+\hat U$ on one
    microscopic state space, and Burt transforms one complete heterogeneous
    microscopic problem into a fixed envelope basis. The present project must
    first identify independently produced pristine and doped retained spaces
    before their operators can be subtracted.
  \item \textbf{General exact kernel versus reduced impurity model.}
    Luttinger--Kohn begin with a multiplicative $U(\mathbf r)$, while Burt's
    exact transformation can retain nonlocal and multicomponent structure. A
    local scalar impurity potential nevertheless emerges only after a separately
    justified localization and model-class reduction.
  \item \textbf{Analytical periodic basis versus Wannier construction.}
    Luttinger--Kohn use band-edge periodic functions, Burt permits one fixed
    complete periodic basis or a wave-vector-dependent Wannier--Slater basis,
    and Ermoneit et al. choose valley sectors of a strained-silicon Bloch basis.
    None of these operations performs modern composite-band disentanglement,
    maximally localized gauge selection, or independent basis alignment.
  \item \textbf{Asymptotic or Fourier scale separation versus finite
    computational evidence.} Luttinger--Kohn assume a gentle perturbation and
    small relevant wave vectors; Burt tests whether the envelopes are globally
    dominated by retained Fourier sectors; Ermoneit et al. enforce disjoint
    valley support and test constant-shift invariance. The project must
    additionally test finite supercells, reciprocal meshes, retained ranks,
    embeddings, and spatial domains before making a crossover claim.
  \item \textbf{Reduction fidelity versus parent-model validity.} Agreement
    with a Luttinger--Kohn-type continuum equation can establish reduction
    fidelity to a selected Kohn--Sham parent. It does not by itself validate the
    parent against experiment.
\end{enumerate}

\section{Implications for phosphorus and boron in silicon}
\index{phosphorus!Luttinger--Kohn interpretation}
\index{boron!Luttinger--Kohn interpretation}

For phosphorus, Luttinger--Kohn support beginning with a multivalley envelope
space associated with the equivalent silicon conduction minima. They also support
separating the smooth long-range regime from the short-range region where the
perturbation varies strongly within a cell. It does not determine the
intervalley central-cell operator, the correct valley combinations, or the
mapping from an independently generated doped calculation into the pristine
band-edge basis.

For boron, Luttinger--Kohn support a coupled-envelope description of a degenerate
valence edge and warn against a scalar single-band treatment when spin--orbit
splitting is not large compared with the impurity scale. They do not select the
modern parent calculation, spinor basis, continuum parameters, or validation
observables required by this project's boron branch.

Burt's fixed-basis extraction adds a testable intermediate question for both
dopants: whether the exact represented envelopes are concentrated in the
retained reciprocal sectors strongly enough to justify the local multiband
operator~\eqref{eq:burt-local-multiband}, and then whether a dominant-envelope
elimination is accurate for the declared observable. Ermoneit et al. sharpen the
phosphorus-side question by supplying explicit sector projectors and a
constant-energy-shift invariance test for intervalley coupling. Their paper
studies Si/SiGe conduction valleys, not substitutional phosphorus, boron, or a
central-cell impurity operator; application to either dopant remains proposed
work.

The two dopants therefore share the need for a declared microscopic-to-envelope
map, but they do not share the same continuum state space. Valley multiplicity
organizes the phosphorus donor problem; valence-band degeneracy and spin--orbit
coupling organize the boron acceptor problem.

\section{Error and evidence interpretation}
\index{error!Luttinger--Kohn approximation}
\index{evidence!literature analysis}

The approximations discussed by Luttinger--Kohn, Burt, and Ermoneit et al.
should be separated into the error categories used by this monograph.

\begin{description}
  \item[Parent-model error.] The periodic one-particle parent may disagree with
    the physical material. The effective-mass derivation does not remove this
    discrepancy.
  \item[Scale-separation and analytical reduction error.] Neglected reciprocal-
    lattice Fourier components, localization of the exact nonlocal envelope
    kernel, omitted interface terms, spectral leakage between valley sectors,
    higher powers of wave vector, remote-envelope elimination, residual
    interband mixing, and short-range failure contribute
    to the difference between the microscopic parent and the envelope model.
  \item[Numerical error.] Any finite representation of the parent bands,
    derivatives, envelopes, or coupled PDE adds discretization and solver
    errors not analyzed merely by the historical formal derivation.
  \item[Model-class error.] A scalar, multivalley, or multiband continuum class
    may omit operator structure required by the aligned atomistic target.
\end{description}

The present reading establishes only what the cited papers state and how their
assumptions relate to the declared project question. It is literature analysis,
not a calculated result, numerical verification, scientific validation, or
proof of the project's proposed continuum limit.

\section{Summary of what the papers establish}
\index{Luttinger--Kohn model!summary}
\index{Burt envelope-function theory!summary}
\index{Ermoneit multivalley theory!summary}

Luttinger and Kohn go beyond the curvature mnemonic by expanding the perturbed
state in a complete orthonormal family built from band-edge periodic functions
and by using a canonical transformation to remove interband coupling through
the retained quadratic order~\cite{luttingerKohn1955}. For a simple band and a
potential that varies little across a unit cell, their momentum-space equation
(II.35) becomes the coordinate-space envelope equation (II.38), and the leading
microscopic state has the product form in (II.40). Their derivation treats
minima away from the Brillouin-zone center and multiple equivalent valleys in
Section~III, degenerate bands through coupled envelope equations in Section~IV,
and spin--orbit coupling in Section~V.

These are approximation statements. The relevant envelope or impurity scale
must be large compared with the lattice spacing, the perturbation must be
gentle away from any explicitly short-range region, and the band expansion is
retained only through quadratic order. Lowest-order multivalley theory does not
determine the correct valley combinations, and a single-band description is
not justified when coupling to a nearby degenerate or spin--orbit-split manifold
is comparable with the perturbation scale.

Burt makes the envelope transformation itself exact by equations (3.1) and
(3.8) of the 1992 paper and equations (2.4)--(2.5) of the 1999 paper. The local
multiband equation (1992, equation (5.1)), dominant-envelope elimination
(1992, equations (6.3)--(6.10)), interface conditions (1992, equations
(7.1)--(7.4)), operator ordering (1999, equations (4.4)--(4.11)), and
out-of-zone reconstruction (1999, equations (5.1)--(5.11)) are subsequent
approximations or consequences with separately identifiable assumptions. The
exact representation does not make a truncated local effective-mass model
exact.

Ermoneit et al. make the multivalley decomposition unique through equations
(5) and (9)--(11), retain the project--multiply--project kernel (14) in the
single-band equation (19), and prove the constant-shift identity (26). Their
first-order intervalley matrix element (30) gives
$E_{\mathrm{VS}}=2|\Delta|$, while equations (39)--(40) define a filtered-local
approximation that restores the same constant-shift invariance. This is a
concrete extension of Burt's exact representation logic, not evidence that an
unprojected local multivalley model is exact.

\documentcallout{Source review resolved}{Luttinger--Kohn (1955), key
\texttt{luttingerKohn1955}~\cite{luttingerKohn1955}, Burt (1992, 1999), keys
\texttt{burt1992}~\cite{burt1992} and \texttt{burt1999}~\cite{burt1999}, and
Ermoneit et al. (2026), key \texttt{ermoneit2026}~\cite{ermoneit2026}, have
been read in full. Short Burt references: M.~G. Burt,
\emph{J. Phys.: Condens. Matter} \textbf{4}, 6651--6690 (1992),
\href{https://doi.org/10.1088/0953-8984/4/32/003}{doi:10.1088/0953-8984/4/32/003};
and M.~G. Burt, \emph{J. Phys.: Condens. Matter} \textbf{11}, R53--R83 (1999),
\href{https://doi.org/10.1088/0953-8984/11/9/002}{doi:10.1088/0953-8984/11/9/002}.
Ermoneit reference: L.~Ermoneit, A.~Thayil, T.~Koprucki, and M.~Kantner,
\emph{Phys. Rev. B} \textbf{113}, 245306 (2026),
\href{https://doi.org/10.1103/md2x-s44y}{doi:10.1103/md2x-s44y}.
These sources do not establish modern Wannier disentanglement, independently
generated pristine--doped alignment, or the project's proposed continuum
crossover.}

\section{Questions carried forward}

The papers leave the present project with five concrete questions:
\begin{enumerate}
  \item What explicit identification maps an accepted retained Kohn--Sham space
    to the donor or acceptor envelope space?
  \item Over what tested spatial domain is the aligned impurity operator gentle
    relative to the lattice scale?
  \item Which intervalley or interband terms remain after the declared lattice
    reduction?
  \item Which operator norm and observable tests establish the finite crossover
    from the atomistic model to the continuum class?
  \item How do parent-model, numerical, and reduction discrepancies remain
    separate when the continuum model is compared with experiment?
\end{enumerate}
Answering these questions requires the modern representation, alignment, and
evidence program developed in the main chapters. The four papers supply the
foundational representation, reduction pattern, and multivalley support
diagnostic but not completed answers.
